\subsection{補助利得に対する定量評価}
\label{sec:ds-estimates}
第\ref{sec:components}節で使う数値評価を示す。
以下の成分は一つの確率 $Q$ の下での零始点正準成分である。
取引の帰還先の価格は常に $S$ とする。構成で使う操作は有界予測可能制限、加法、
正のスカラー倍、閉停止だけである。系\ref{cor:market-operations}によって、構成した許容利得は
全て元市場に実現される。議論は\cite[Lemmas 4.7--4.10]{DS}に従うが、
実現済み利得の族として述べるため、補助分解は添字ごとに異なってよい。

\begin{lemma}[具体的なjump包絡]\label{lem:explicit-jump-envelope}
$Y=M+A$ を予測可能FV成分を持つ正則special過程とし、
$0\leq q\in L^2(Q)$ により全時刻で $|Y_t|\leq q$ とする。
各有限 $T$ に対して非負の $\xi_T\in L^2(Q)$ が存在し、
\[
 |\Delta M_t|\leq\xi_T\quad(0<t\leq T),\qquad
 \|\xi_T\|_2\leq6\|q\|_2
\]
となる。
\end{lemma}
\begin{proof}
$Z_t=E_Q[q\mid\mathcal F_t]$ のcàdlàgマルチンゲール版を取り、
$Z_T^*=\sup_{t\leq T}|Z_t|$ と置く。
予測可能停止時刻 $\sigma\leq T$ では、正準分解のjumpの条件付き期待値から
\[
 \Delta A_\sigma=E_Q[\Delta Y_\sigma\mid\mathcal F_{\sigma-}],\qquad
 |\Delta A_\sigma|\leq2Z_{\sigma-}
\]
を得る。局所成分については先に局所化し、$\sigma$ を下から予告する。
$Y$ の可積分包絡により、条件付き期待値をjumpの極限へ移せる。
予測可能過程のjumpの列挙によって、$A$ の非零jumpは可算個のそのような時刻で覆われる。
共通の確率1の集合で全ての $0<t\leq T$ に対して $|\Delta A_t|\leq2Z_T^*$ となる。
従って
\[
 \xi_T=2q+2Z_T^*,\qquad
 \|\xi_T\|_2\leq2\|q\|_2+4\|q\|_2=6\|q\|_2
\]
と取れる。最後はDoobの $L^2$ 不等式である。
これは\cite[Corollary 2.4]{DS}のjump評価について、包絡とhorizonを明示したものである。
\end{proof}

\begin{lemma}[有限horizonの消失リスク戦略]\label{lem:quantitative-47}
元市場に実現された利得 $Y=M+A$ が $Y\geq-1$、$|Y|\leq q$、
$\|q\|_2\leq B$ を満たすとする。
$0<\alpha\leq1/4$、正整数 $n$、有限 $T>0$ に対し
\[
 Q\left(\sup_{t\leq T}|M_t|>n^3\right)>8\alpha
\]
とする。$k=\lfloor\alpha n/4\rfloor$ と置き、$k>0$、
$6B\leq n^2$、$(B/n)^2\leq\alpha$、
$m_n=\alpha^2-(4B/(\alpha n^2))^2\geq0$ と仮定する。
元市場の利得 $V_n$ で、$T$ 以後一定であり、
\[
 \begin{aligned}
 V_n(t)&\geq-n^{-1/4}-\frac{2}{nk}\quad(\forall t),\\
 Q\left(V_n(T)\geq\frac{\alpha m_n}{4}-n^{-1/4}\right)
 &\geq\frac{m_n}{2}-\frac{16\sqrt n}{k}
 \end{aligned}
\]
を満たすものが存在する。
\end{lemma}
\begin{proof}
最初の同時passageと再尺度化した利得を
\[
 \theta=T\wedge\inf\{t:|M_t|>n^3\text{ または }|Y_t|>n\},\qquad
 L=n^{-2}Y^\theta=N+B^L
\]
とする。停止前の利得上界と $Y\geq-1$ から、$\theta$ でのjumpを含めて
$\Delta L\geq-(n+1)/n^2\geq-2/n$ である。また $L^*\leq q/n^2$ である。
前補題により $N$ のjump包絡の $L^2$ ノルムは $6B/n^2\leq1$ 以下となる。
停止前には $|N|\leq n$ であるから、全時間の包絡は $n$ とこのjump包絡の和で抑えられる。
従って $N$ は真の $L^2$ マルチンゲールである。
Chebyshev不等式で利得上限に達する事象を除き、初回passageでの右連続性を使うと
\[
 Q(N_T^*>n/2)>8\alpha-(B/n)^2\geq7\alpha
\]
を得る。半径を $n/2$ とすることで、passageでの厳密不等号の問題を避ける。

$\kappa_0=0$ から始め、$\kappa_{j-1}$ 以後の $N$ の変位の絶対値が
初めて $1$ を厳密に超える時刻と $T$ の最小を $\kappa_j$ とする。
$1\leq j\leq k$ に対し $F_j=N_{\kappa_j}-N_{\kappa_{j-1}}$ と置く。
jump包絡より $\|F_j\|_2\leq2$、任意抽出定理より $E_QF_j=0$ である。
最初の $k$ 個のexcursionが完了していなければ、最大関数はこれらの停止excursionの
最大関数の和で抑えられる。Doob不等式はそれぞれの $L^2$ ノルムを $4$ で抑えるので、
\[
 Q(N_T^*>n/2,\ k\text{ 個のexcursionが未完了})
 \leq\left(\frac{4k}{n/2}\right)^2\leq\alpha.
\]
特に全ての $j\leq k$ に対し $Q(|F_j|\geq1)>6\alpha$ である。
従って $EF_j^-=E|F_j|/2>3\alpha$ となる。
$E_j=\{F_j^->\alpha\}$ とすると $E[F_j^-1_{E_j}]>2\alpha$ であり、
Cauchy--Schwarz不等式で同じ期待値は $2\sqrt{Q(E_j)}$ 以下である。
従って $Q(E_j)>\alpha^2$ となる。
対応する利得増分の $L^2$ ノルムは $2B/n^2$ 以下である。
水準 $\alpha/2$ のChebyshev不等式と $B^L=L-N$ から
\[
 Q(B^L_{\kappa_j}-B^L_{\kappa_{j-1}}\geq\alpha/2)\geq m_n
\]
を得る。

$dB^L$ の予測可能な正のHahn集合 $P$ を一つ選び、実際の利得を
\[
 R=(1_P1_{(0,\kappa_k]})\cdot L=N^R+C
\]
とする。$C$ は非減少で、各excursionでの増分は $B^L$ の対応する増分以上である。
上で得た $k$ 個の事象の成立個数を $Z$ とすると、$0\leq Z\leq k$、$EZ\geq km_n$ である。
$EZ\leq km_n/2+kQ(Z\geq km_n/2)$ より
\[
 Q(C_T\geq k\alpha m_n/4)\geq m_n/2
\]
となる。excursion事象の独立性は使わない。
予測可能制限は各マルチンゲールexcursionの $L^2$ ノルムを縮小する。
互いに交わらない停止区間の増分の直交性から $E|N^R_T|^2\leq4k$、従って
\[
 Q\left(\sup_{t\leq T}|N^R_t|>d\right)\leq16k/d^2
\]
を得る。ここでは全て固定horizonの真の $L^2$ マルチンゲールであり、
任意抽出と等長性がノルム縮小と直交性を正当化する。

最後に
\[
 d=kn^{-1/4},\qquad
 \lambda=T\wedge\inf\{t:R_t<-d\},\qquad V_n=k^{-1}R^\lambda
\]
と置く。正のHahn集合への制限は利得のjumpまたは零を選ぶから、$\Delta R\geq-2/n$ である。
従って停止時の飛び越しも含めて $R^\lambda\geq-d-2/n$ となる。
上のマルチンゲール最大事象の外では、$T$ まで $R=N^R+C\geq-d$、
かつ $R_T\geq C_T-d$ である。$C_T$ の事象と交差を取り $k$ で割れば、
$16k/d^2=16\sqrt n/k$ より両評価を得る。
$L,R,V_n$ を定義する全操作は本小節冒頭の元市場への実現を持ち、
最後の停止によって終端値は文字どおり $V_n(T)$ となる。
\end{proof}
この有限尺度の構成を用いるNFLVRの矛盾は、命題\ref{prop:indexed-maximal}で行う。

\begin{lemma}[凸尾部の許容性と具体的な包絡]\label{lem:quantitative-tails}
各利得が有限時刻以後一定である実現済み族 $Y_i=M_i+A_i\geq-1$ を考える。
共通包絡 $q\geq0$ は $L^2(Q)$ に属し、$\{M_i^*\}_i$ は確率有界と仮定する。
$\tau_i(c)=\inf\{t:|M_i(t)|>c\}$、$E_i(c)=\{\tau_i(c)<\infty\}$ と置く。
有限凸重み $w$ に対する尾部利得とそのM成分を
\[
 Z^{w,c}=\sum_iw_i(Y_i-Y_i^{\tau_i(c)}),\qquad
 L^{w,c}=\sum_iw_i(M_i-M_i^{\tau_i(c)})
\]
とする。利得の包絡は
\[
 G^{w,c}=2q\sum_iw_i1_{E_i(c)},\qquad
 \sup_w\|G^{w,c}\|_2\leq2\sup_i\|q1_{E_i(c)}\|_2\longrightarrow0
\]
で与えられる。各有限horizonでjump包絡 $\xi_T^{w,c}$ を持ち、
$\sup_w\|\xi_T^{w,c}\|_2\leq6\sup_w\|G^{w,c}\|_2\to0$ となる。
このノルム上界のcutoffは $T$ に依存しない。
さらに任意の $\varepsilon>0$ に対し $N\geq0$ を選べる。
各 $0<\delta\leq\varepsilon/4$ に対して、重みを選ぶ前に $c_0$ を定められ、
$c\geq c_0$ かつ $L^{w,c}$ の水準 $\varepsilon$ への有限passageの確率が
$\varepsilon$ を超えるなら、元市場の利得 $V=M^V+A^V$ と包絡 $g$ が存在して
\[
 V\geq-1,\qquad |V|\leq g,\qquad \|g\|_2\leq1,\qquad
 Q\!\left((M^V)^*>\frac{\varepsilon}{2(N+1)\delta}\right)>\frac\varepsilon2
\]
となる。この主張は一つの有限凸尾部の構成であり、NFLVRを仮定しない。

\end{lemma}
\begin{proof}
仮定した最大関数の確率有界性から $\sup_iQ(E_i(c))\to0$ である
（必要ならpassage水準を少し下げて適用する）。
$q^2$ の積分の絶対連続性により $\sup_i\|q1_{E_i(c)}\|_2\to0$ となる。
各尾部はpassage以前およびpassage時に零で、それ以後の絶対値は $2q$ 以下である。
包絡とそのノルム評価は三角不等式から従う。

有限尺度の構成を示す。$\varepsilon>0$ を固定する。
$Q(q\geq N)<\varepsilon/4$ となる $N$ を選び、全利得の絶対値のいずれかが
初めて $N$ を厳密に超える時刻を $\zeta$ とする。
$0<\delta\leq\varepsilon/4$ に対し、全ての $i$ で $Q(E_i(c))\leq\delta^2$、
全ての重みで $\|G^{w,c}\|_2\leq\delta$ となるcutoffを選ぶ。
これらは結論を破る重みを選ぶ前の選択である。
active mass
\[
 F_t=\sum_iw_i1_{\{\tau_i(c)<t\}},\qquad
 \gamma=\inf\{t:F_t>\delta\}
\]
は非減少、適合、左連続であるから $F_\gamma\leq\delta$ となる。
一方 $EF_\infty\leq\delta^2$ より $Q(\gamma<\infty)\leq\delta$ である。
$t\leq\zeta\wedge\gamma$ なら全てのactiveな開始時刻は $\zeta$ より厳密に前であるため、
$Y_i(\tau_i(c))\leq N$ である。また閉停止時を含めて $Y_i(t)\geq-1$ なので
\[
 Z^{w,c}_t\geq-(N+1)F_t\geq-(N+1)\delta
\]
となる。$L^{w,c}$ の水準 $\varepsilon$ へのpassageも停止に加える。
このpassageが有限である確率が $\varepsilon$ を超えるとする。
利得上限とactive massによる二つの停止が先に起こる事象を除いても、残る確率は
$\varepsilon-\varepsilon/4-\delta\geq\varepsilon/2$ より大きい。
その事象では閉停止時のM成分の絶対値が $\varepsilon$ 以上であり、
特に $\varepsilon/2$ を厳密に超える。有限和に現れる各利得は有限時刻後一定なので、
この停止利得も有限時刻後一定である。
$(N+1)\delta$ で正規化すれば、元市場の $1$ 許容利得を得る。
包絡ノルムは $1/(N+1)\leq1$ であり、表示した増幅後の確率下界が従う。
最後に補題\ref{lem:explicit-jump-envelope}を $q=G^{w,c}$ として適用する。
具体的には
\[
 \xi_T^{w,c}=2G^{w,c}+2\sup_{t\leq T}|E_Q[G^{w,c}\mid\mathcal F_t]|
\]
と取れる。ノルム評価はDoob不等式から従い、horizonに依存しない。
\end{proof}

\contractref{F1}は補題\ref{lem:quantitative-47}の終端 $f=V_n(T)$ を取り、
$K_0\subseteq K_0-L^0_+$ を使ったものである。\contractref{F2}は補題\ref{lem:quantitative-tails}の結論そのものである。
